\documentclass[12pt]{article}
\usepackage[T1]{fontenc}
\usepackage[utf8]{inputenc}
\usepackage{lmodern}
\usepackage{textcomp}
\usepackage{enumitem}
\usepackage{geometry}
\geometry{margin=1in}
\begin{document}
\begin{center}
Answer Key - Constitutional Law - Graduate Level Assessment 7 \
\small (Answers are ordered exactly as in the question paper.)
\end{center}

\section*{Multiple Choice}
\begin{enumerate}[label=\arabic*.]
\item \textbf{Answer:} A
\\ \textit{Options:} A) not hearsay.; B) hearsay, but admissible as an admission.; C) hearsay, but admissible as a dying declaration.; D) hearsay not within any recognized exception.
\\ \textit{Note:} Correct option: A - not hearsay.
\item \textbf{Answer:} D
\\ \textit{Options:} A) the causal relationship between the defendant's act and the resulting death.; B) the attendant circumstances surrounding the death.; C) the nature of the act causing the death.; D) the defendant's state of mind at the time the killing was committed.
\\ \textit{Note:} Correct option: D - the defendant's state of mind at the time the killing was committed.
\item \textbf{Answer:} B
\\ \textit{Options:} A) upon court order.; B) upon request.; C) upon request and showing of good cause.; D) under no circumstances.
\\ \textit{Note:} Correct option: B - upon request.
\item \textbf{Answer:} D
\\ \textit{Options:} A) interpretation of treaties; B) interpretation of admiralty laws; C) disputes between states and foreign citizens; D) all of the above
\\ \textit{Note:} Correct option: D - all of the above
\item \textbf{Answer:} A
\\ \textit{Options:} A) It is an agreement between two or more to commit a crime.; B) It is solicitation of another to commit a felony.; C) No act is needed other than the solicitation.; D) Withdrawal is generally not a defense.
\\ \textit{Note:} Correct option: A - It is an agreement between two or more to commit a crime.
\end{enumerate}

\section*{True / False}
\begin{enumerate}[label=\arabic*.]
\setcounter{enumi}{5}
\item \textbf{Answer:} True
\\ \textit{Options:} A) True; B) False
\\ \textit{Note:} LLM-generated true statement based on MCQ.
\item \textbf{Answer:} True
\\ \textit{Options:} A) True; B) False
\\ \textit{Note:} LLM-generated true statement based on MCQ.
\item \textbf{Answer:} False
\\ \textit{Options:} A) True; B) False
\\ \textit{Note:} LLM-generated false statement. Correct answer: 'Is responsible for damages caused by third-party tortfeasor.'.
\item \textbf{Answer:} True
\\ \textit{Options:} A) True; B) False
\\ \textit{Note:} LLM-generated true statement based on MCQ.
\item \textbf{Answer:} False
\\ \textit{Options:} A) True; B) False
\\ \textit{Note:} LLM-generated false statement. Correct answer: 'Crops and timber to be severed from the property next summer.'.
\end{enumerate}

\section*{Long Form Response}
\begin{enumerate}[label=\arabic*.]
\setcounter{enumi}{10}
\item \textbf{Answer:} A three-fourths vote in the Senate is necessary to convict and remove from office. Explain why this is the correct answer and provide supporting evidence. Reference key facts or concepts that support this choice.
\\ \textit{Note:} Long-form derived from MCQ; award credit for stating the correct idea and giving supporting reasoning.
\item \textbf{Answer:} The defendant is compelled to stand trial in street clothing.. Explain why this is the correct answer and provide supporting evidence. Reference key facts or concepts that support this choice.
\\ \textit{Note:} Long-form derived from MCQ; award credit for stating the correct idea and giving supporting reasoning.
\end{enumerate}

\vfill
\noindent\textit{End of Answer Key}
\end{document}